\documentclass[11pt]{article}
\usepackage{ochamath}
\subjectcolor{2F5DA8}
\setsheet{線形代数}{線形写像と基底変換　演習}{https://university-math-crj.pages.dev/linear-algebra/linear-maps/}
\begin{document}
\makesheethead{解答}

\problem[線形性]
\begin{ans}
$T$ は行列 $\begin{pmatrix}2&-1\\0&3\end{pmatrix}$ で表せるので線形。$S(0,0)=(1,0)$ は零でないため、$S$ は線形ではない。
\end{ans}
\point{非線形の判定には1つの反例で足りる。}

\problem[核と像]
\begin{ans}
核の基底は $(-1,-1,1)^{\mathsf T}$。像は $\mathbb R^2$ 全体で標準基底 $(1,0)^{\mathsf T},(0,1)^{\mathsf T}$ を取れる。核1次元＋像2次元＝入力3次元。
\end{ans}
\point{出力空間の次元と像の次元が同じなら全射。}
\newpage

\problem[同じ基底での表現]
\begin{ans}
$P^{-1}=\begin{pmatrix}1&-1\\0&1\end{pmatrix}$。$A(1,0)^{\mathsf T}=2(1,0)^{\mathsf T}$、$A(1,1)^{\mathsf T}=5(1,1)^{\mathsf T}$ より、$P^{-1}AP=\operatorname{diag}(2,5)$。検算：$AP=\begin{pmatrix}2&5\\0&5\end{pmatrix}=P\operatorname{diag}(2,5)$。
\end{ans}
\point{表現行列の列は基底ベクトルの像を基底座標で表したもの。}

\problem[異なる基底]
\begin{ans}
$Q^{-1}AP=\operatorname{diag}(2,1/3)$。出力基底座標は $(2,1)^{\mathsf T}$。標準座標では入力は $(2,3)^{\mathsf T}$、出力も $Q(2,1)^{\mathsf T}=(2,3)^{\mathsf T}$ で同じ。
\end{ans}
\point{恒等写像でも入力と出力の基底が違えば表現行列はIとは限らない。}
\end{document}
